\chapter{Negation and Negative Concord}\label{ch:negation}

\section{Polarity domains and their exponents}\label{sec:clausal-negation}

Indicative verbal predicates take \ird{za-}/\ird{z-}, zero-copula
predicates take \ird{staly}, and existential predicates alternate
\ird{jec-} with \ird{zad-}; non-indicative clauses instead realise
negative polarity through their dedicated Mood suffixes. These forms express
one local negation, including when negative indefinites participate in concord.
The distribution is summarised in Table~\ref{tab:neg-exponents}.

\begin{table}[ht]
  \footnotesize
  \caption{Negation strategies by predicate/mood type.}\label{tab:neg-exponents}
  \begin{tabularx}{\textwidth}{L L L L}
    \toprule
    \textsc{predicate/mood type} & \textsc{affirmative} & \textsc{negative} & \textsc{interpretation} \\
    \midrule
    Verbal (indicative)           & verbal complex                & \ird{za-} prefix; \synhead{NegP} projected          & propositional \\
    Equative / zero copula        & $\varnothing$ copula            & \ird{staly} (free morpheme, clause-final) & predicational \\
    Identificational (part.\ nom.) & part.\ nom.\ + $\varnothing$  & part.\ nom.\ + \ird{staly}                & identificational \\
    Existential                   & \ird{jec-}                    & \ird{zad-} (suppletive root); zero exponent of Neg\textsuperscript{0} & propositional \\
    \midrule
    \multicolumn{4}{l}{\textit{Negative Mood without a separate NegP\rlap{\textsuperscript{†}}}} \\[2pt]
    Command (imperative)          & \ird{-ím}                     & \ird{-é} (prohibitive)                    & directive non-performance \\
    Suggestion (hortative)        & \ird{-ká}                     & \ird{-ku} (neg.\ hortative)               & directive non-performance \\
    Hypothetical (irrealis)       & \ird{-ěv}                     & \ird{-ir} (neg.\ irrealis)                & hypothetical non-occurrence \\
    \bottomrule
    \multicolumn{4}{p{\dimexpr\linewidth-2\tabcolsep\relax}}{\textsuperscript{†}\footnotesize
      Negative mood suffixes are analysed as exponents of negative Mood, with
      the prohibitive retaining directive force over a negated proposition.
      Negative irrealis composes locally as $\mathrm{IRR}(\neg P)$. See
      §\,\ref{sec:vm-moodp}.} \\

  \end{tabularx}
\end{table}

\section{Indicative verbal negation}\label{sec:negation}

\subsection{Forms, distribution, and prosody}\label{sec:neg-allomorphy}

Indicative verbal negation is expressed by unstressed \ird{za-}, which attaches
as a proclitic to the left edge of the verbal complex without displacing the
host's primary stress. The imperative, hortative, and irrealis have the negative
mood suffixes \ird{-é}, \ird{-ku}, and \ird{-ir}, respectively
(§\,\ref{sec:mood}). Nominal predicates take clause-final \ird{staly}, while
existential predicates distinguish affirmative \ird{jec-} from negative
\ird{zad-} (§\,\ref{sec:clausal-negation}).

The negative marker is underlyingly /za/, with its short vowel retained before
a consonant, as in \ird{za-} + \ird{pjaštník} $\rightarrow$
\ird{zapjaštník}, and deleted before a vowel, giving \ird{z-}. Deletion applies before both short and long vowels, irrespective of stress,
and is morphologically restricted to the negative marker. The choice of form consequently depends on the host's
phonological onset: \ird{i} always represents a vowel, whereas the labial
digraph \ird{pj} represents a consonant (§\,\ref{sec:orthographic-i}).

\begin{table}[htbp]
  \centering
  \caption{Indicative negative allomorphy; vowel and glide sequences are schematic.}\label{tab:neg-allomorphy}
  \begin{tabularx}{\textwidth}{l l X}
    \toprule
    {\sc host onset} & {\sc spelling} & {\sc analysis} \\
    \midrule
    Consonant & \ird{za-} & Short /a/ retained; e.g.\ \ird{zapjaštník}. \\
    \ird{a, á} & \ird{za\ldots, zá\ldots} & \ird{z-} + host vowel; length belongs to the host. \\
    \ird{e, é} & \ird{ze\ldots, zé\ldots} & \ird{z-} + unchanged host vowel. \\
    \ird{i, í} & \ird{zi\ldots, zí\ldots} & \ird{z-} + unchanged host vowel; /z/ remains alveolar. \\
    \ird{o, ó} & \ird{zo\ldots, zó\ldots} & \ird{z-} + unchanged host vowel. \\
    \ird{u, ú} & \ird{zu\ldots, zú\ldots} & \ird{z-} + unchanged host vowel. \\
    \ird{y, ý} & \ird{zy\ldots, zý\ldots} & \ird{z-} + unchanged host vowel. \\
    \ird{j} /j/ & \ird{zaj\ldots} & Consonantal onset; full \ird{za-} retained. \\
    \bottomrule
  \end{tabularx}
\end{table}

Before \ird{i, í}, the reduced marker retains alveolar /z/: the high front
vowel does not select a postalveolar consonant
(§\,\ref{sec:sibilant-palatalisation}). The glide /j/, being consonantal,
takes \ird{za-}, so that a \ird{j}-initial host begins \ird{zaj\ldots}
under negation. The negative marker has no \ird{ž-} allomorph.

In weak position, the marker's short /a/ may centralise to [ɐ], giving
[za-] $\sim$ [zɐ-] without any change to the host's stress
(§\,\ref{sec:stress}). Since the marker has no long-vowel variant, an initial
\ird{zá\ldots} contains \ird{z-} followed by the host's own long \ird{á}.
Throughout this grammar, the label \ird{za-} includes its regular \ird{z-}
allomorph.

\subsection{Formal analysis: \synhead{NegP} and the position of the negative marker}
\label{sec:vm-neg}
\label{sec:vm-neg-exponents}

Indicative clausal negation is analysed as a \synhead{NegP} above
\synhead{AspP} and below \synhead{MoodP}
(Table~\ref{tab:th-clause-spine}), whose realisation depends on the predicate.
In ordinary verbal clauses, \ird{za-} precedes the root and forms a prosodic
unit with the inflected verb, while in zero-copula clauses \ird{staly} follows
the nominal predicate. With existential predicates, negative polarity is
expressed by suppletion of the root, \ird{jec-}/\ird{zad-}.

Iridian exhibits \term{strict negative concord}: negative indefinites such as
\ird{deno} \trsl{nothing} and \ird{zide} \trsl{no one} require the
negative realisation appropriate to their finite clause, while their
co-occurrence expresses a single semantic negation
(§\,\ref{sec:nm-negative-concord}). The grammar analyses this dependency through
a silent operator $\mathrm{Op}_{\neg}$ bearing an interpretable negative
feature $[i\textsc{neg}]$, together with an uninterpretable feature
$[u\textsc{neg}]$ on Neg\textsuperscript{0} and on each concord associate.
The operator occupies Spec,\synhead{NegP}; it supplies semantic negation,
whereas the feature on Neg\textsuperscript{0} determines its participation in
the dependency and its morphological realisation. This division of labour
follows the syntactic approach to negative dependencies developed by
Zeijlstra and Penka \cite{zeijlstra2022,penka2011}; its application to Iridian is the analysis
adopted here.

The exponent of verbal Neg\textsuperscript{0} has a stipulated left-edge
instruction in the Iridian mapping to PF. Let $H$ be the independently formed
inflected verbal host, including any lexical preverb and licensed valency,
orientation, aspect, mood and evidential exponents. Vocabulary realisation
places \ird{za-} immediately before $H$, after which its regular onset-conditioned
allomorphy and prosodic attachment apply (§\,\ref{sec:neg-allomorphy}):
\[
 \mathrm{Neg}^{0}_{\text{verbal}}+H
 \quad\longmapsto\quad \ird{za-}H.
\]
The plus sign represents the realisation dependency, not surface precedence.
The instruction targets this verbal host, excluding arguments and free
particles; it preserves the host's internal constituency and suffix order.
Neg remains above Asp in the interpreted structure, although its exponent
precedes the complete verbal word.

Affix-specific directionality is among the interface mechanisms discussed by
Harley \cite{harley2011}; the choice of a verbal host and its left edge is an
Iridian-specific realisation rule. No syntactic head movement or sequence of
post-syntactic swaps derives this position. In particular, the rule does not
apply Local Dislocation across intervening material: that operation is
restricted by linear adjacency in Embick and Noyer's account
\cite{embicknoyer2007}. The same distinction between hierarchical position and
an explicit PF precedence instruction governs speaker-anchoring \ird{si}
(§\,\ref{sec:nm-si-linearisation}). Prefixal negation therefore does not
independently establish either Neg's height or an unrestricted exception to
the Mirror Principle.

In zero-copula clauses such as \ird{Marek doktor staly} \trsl{Marek is not a
doctor}, invariant \ird{staly} realises a head-final Neg selecting the nominal
predication directly, without an intervening \synhead{AspP}. In verbal
clauses, Neg instead selects an aspectual predicate, including a finite
stative predicate (§\,\ref{sec:vm-tension-stative}). The difference in
complement type corresponds to the distribution of \ird{staly} and
\ird{za-}.

Existential negation is expressed by \ird{zad-}, the negative allomorph of
Exist\textsuperscript{0}, with a zero exponent at Neg\textsuperscript{0}.
Thus \ird{za-} and \ird{staly} spell out Neg itself, whereas \ird{zad-}
realises the existential root in the environment of Neg. Root suppletion is
well documented among negative existentials \cite{veselinova2016}; in the
Iridian analysis, the negative head conditions the suppletive root while
remaining distinct from it in the syntactic representation.

\begin{figure}[ht]
  \footnotesize
  \caption{Proposed structures for indicative verbal, equative, and existential
  negation. The verbal marker attaches to its host at PF, while existential
  negation combines a zero exponent of Neg with a negative root allomorph.
  Higher topic structure is omitted; Spec,\synhead{NegP} hosts the silent
  negative operator.}\label{fig:neg0-comparative}
  \medskip
  \centering
  \scriptsize
  \begin{minipage}[t]{0.30\textwidth}
    \centering
    \textsc{Verbal clause}\par
    \medskip
    \begin{forest}
      syntax tree, for tree={l sep=8pt,s sep=6pt,inner sep=0.5pt}
      [\synhead{NegP}
        [{$\mathrm{Op}_{\neg}$\\{$[i\textsc{neg}]$}}]
        [{Neg$'$}
          [\synhead{AspP}
            [{Asp$'$}
              [\synhead{VP}
                [{V\\{\footnotesize\textit{root}}}]
              ]
              [{Asp\\{\footnotesize\textit{asp.align}}}]
            ]
          ]
          [{Neg$^{0}$\\{$[u\textsc{neg}]$}\\\ird{za-}\\{\small(PF prefix)}}]
        ]
      ]
    \end{forest}
  \end{minipage}
  \hfill
  \begin{minipage}[t]{0.30\textwidth}
    \centering
    \textsc{Equative clause}\par
    \medskip
    \begin{forest}
      syntax tree, for tree={l sep=8pt,s sep=6pt,inner sep=0.5pt}
      [\synhead{NegP}
        [{$\mathrm{Op}_{\neg}$\\{$[i\textsc{neg}]$}}]
        [{Neg$'$}
          [\synhead{EquP}
            [{\synhead{DP}\\{\footnotesize lower copy}}]
            [{Eq$'$}
              [{\synhead{DP}\\{\footnotesize pred.}}]
              [{Eq$^{0}$\\{$\varnothing$}}]
            ]
          ]
          [{Neg$^{0}$\\{$[u\textsc{neg}]$}\\\ird{staly}}]
        ]
      ]
    \end{forest}
  \end{minipage}
  \hfill
  \begin{minipage}[t]{0.35\textwidth}
    \centering
    \textsc{Existential clause}\par
    \medskip
    \begin{forest}
      syntax tree, for tree={l sep=8pt,s sep=6pt,inner sep=0.5pt}
      [\synhead{NegP}
        [{$\mathrm{Op}_{\neg}$\\{$[i\textsc{neg}]$}}]
        [{Neg$'$}
          [\synhead{AspP}
            [{Asp$'$}
              [ExistP
                [{Exist$'$}
                  [{Exist$^{0}$\\\ird{zad-}\\{\small(suppletive)}}]
                ]
              ]
              [{Asp\\{\footnotesize\textit{asp}}}]
            ]
          ]
          [{Neg$^{0}$\\{$[u\textsc{neg}]$}\\{\small($\varnothing$)}}]
        ]
      ]
    \end{forest}
  \end{minipage}
\end{figure}

In Figure~\ref{fig:neg0-comparative}, Neg follows its complement in all
three structures, with the prefixal position of \ird{za-} determined at PF.
An ordinary linked topic occupies the independently projected
 Spec,\synhead{TopP}, above the structures shown here
 (§\,\ref{sec:topic-comment}). Quantificational pivots can be pronounced in the ordinary placement field,
 with lower occurrences of negative indefinites interpreted within the scope
 of the negative operator; topicless existential pivots likewise retain
 their interpretation below that operator (§\,\ref{sec:existential-constructions}). The nominal
 constituents of an equative may permute with topic and comment, while
 \ird{staly} remains clause-final.

In negative existential clauses, \ird{zad-} and the zero exponent of Neg form
a single pattern of contextual exponence, excluding both \ird{*za-jec-} and
\ird{*za-zad-}. Negative mood clauses have a different structure, in which
negative Mood selects the Asp-bearing predicate domain, including any licensed
phasal or comment operators, and no separate Neg head is merged
(§\,\ref{sec:vm-moodp}).

The clitic analysis predicts adjacency between \ird{za-} and its verbal host,
including the exclusion of intervening particles; its behaviour under ellipsis
is undocumented. The invariance of \ird{staly}, which bears no aspect or mood
inflection, supports its treatment as a functional negative marker. The
prefixal and independent forms belong to two of the positional types
represented in Dryer's survey of negative markers \cite{dryer2013neg}.

\section{Mood and polarity}
\label{sec:vm-moodp}

The mood suffixes follow the aspect-alignment exponent in the verbal complex.
Syntactically, Mood selects the Asp-bearing predicate domain, including
licensed phasal and comment layers when present
(§\,\ref{sec:th-exclusive-phasal}). Indicative clausal negation introduces
\synhead{NegP} above Asp and below those operators, whereas negative
non-indicative clauses have no separate Neg head (§\,\ref{sec:vm-neg}). In the indicative, \synhead{MoodP}
is present but has a null exponent. Non-indicative Mood licenses the coarser
$[\pm\textsc{pfv-domain}]$ opposition in the syntax, together with its
constructional orientation requirements. Post-syntactic impoverishment
restricts the features available for Vocabulary Insertion at Asp; it neither
assigns case nor changes the interpretation already established by the
syntactic configuration (§\,\ref{sec:vm-mood-asp}). \synhead{AspP} is thus
present throughout the finite paradigm, although the licensed aspectual
contrasts and their exponents vary with Mood.

Mood scopes over the aspectual predicate and conditions the reduced
Asp.Align paradigm. Its affirmative and negative suffixes are competing
Vocabulary Items at the \synhead{MoodP} head:

\begin{center}
\begin{tabular}{l l l}
\toprule
\textsc{mood} & \textsc{affirmative vi} & \textsc{negative vi} \\
\midrule
Indicative   & $\varnothing$      & ---\quad (\synhead{NegP} projected; \ird{za-}) \\
Imperative   & \ird{-ím}        & \ird{-é} \hspace{3pt}(prohibitive) \\
Hortative    & \ird{-ká}        & \ird{-ku} \\
Irrealis     & \ird{-ěv}        & \ird{-ir} \\
\bottomrule
\end{tabular}
\end{center}

\medskip
\noindent Negative non-indicative clauses are analysed with a Mood head bearing
both a mood specification and $[u\textsc{neg}]$, realised by \ird{-é},
\ird{-ku}, or \ird{-ir}. Negative Mood selects the Asp-bearing predicate domain through any licensed
phasal and comment layers, whereas indicative clausal negation contains a
separate Neg head below those operators and Mood.
Vocabulary Insertion supplies the exponents of these feature bundles after
the syntactic structure has been assembled \cite{embicknoyer2007}. The
absence of a separate Neg node in negative non-indicative clauses accounts for
the exclusion of \ird{za-}.

The interpretable contribution belongs to a silent
$\mathrm{Op}_{\neg}[i\textsc{neg}]$ adjoined to the selected AspP, at the
lower edge of the finite Mood domain. This is a polarity operator, not an
additional functional Neg head or a second Asp head. Its denotation is
$\lambda p.\neg p$; it applies to the proposition abstracted from AspP
before any phasal operator applies. Negative Mood licenses exactly one such
operator in its own finite complement; the dependency cannot cross another
finite Mood or enter an AspP-selecting reduced nominal. Intervening licensed
Phas and comment operators do not interrupt this dependency. Mood's
$[u\textsc{neg}]$ and negative indefinites are concord associates of this
single interpretable operator (§\,\ref{sec:vm-negative-concord}). Thus the
negative suffix is pronounced high, while negation is interpreted low:
\[
 [\mathrm{Mood}_{M,u\mathrm{neg}}\ [C\ [\mathrm{PHAS}\
   [\mathrm{Op}_{\neg}\ [\mathrm{AspP}\ P]]]]]
 \quad\leadsto\quad M(C(\mathrm{PHAS}(\neg P))).
\]
Here $C$ is an independently licensed comment contribution and is omitted
when absent; the specifically licensed lower exclusive composition is defined
in §\,\ref{sec:th-exclusive-phasal}. Without Phas or $C$, the same rule gives
$M(\neg P)$. This is the explicit Iridian analysis of the dependency, rather
than an inference from suffix order. It preserves Mood's selection of an
Asp-bearing domain without requiring adjacency to AspP, and excludes an
additional indicative \ird{za-} in this construction.

The prohibitive \ird{Trava pjašté} \trsl{Do not eat the bread} directs the
addressee to refrain from eating, combining directive force with negative
propositional content \cite{han2001}. The negative hortative similarly
proposes restraint. For negative irrealis, this grammar fixes the local
composition as $\mathrm{IRR}(\neg P)$: non-eating in the relevant hypothetical
situations, rather than denial that eating is possible. This convention does
not settle interactions with independently expressed modal predicates.
Matched phasal and focus diagnostics appear in
§\,\ref{sec:th-polarity-diagnostics}.

The negative irrealis \ird{-ir} triggers \ird{k $\rightarrow$ c} in
\ird{k}-final stems, without the \ird{č}-output or vowel raising associated
with affirmative \ird{-ěv}. Thus \ird{zděk-} \trsl{bring} forms
\ird{zděc-ir}, while \ird{pjašt-} forms \ird{pjaštir}; the forms
\ird{*zděkir} and \ird{*zďíčir} are excluded. The restricted velar-front
grade distinguishes \ird{-ir} from mixed-grade \ird{-ěv} and plain-grade
\ird{-ic}.

\section{Equative negation with \ird{staly}}\label{sec:negative-equatives}

In a zero-copula equative, negation is expressed by clause-final \ird{staly},
whereas the verbal proclitic \ird{za-} requires a verbal host and does not occur
in this construction:

\pex
\a
\begingl
\gla Marek doktor staly.//
\glb Marek.\Dir{} doctor.\Dir{} \Neg{}//
\glft \trsl{Marek is not a doctor.}//
\endgl
\a
\begingl
\gla Doktor Marek staly.//
\glb doctor.\Dir{} Marek.\Dir{} \Neg{}//
\glft \trsl{Marek is not the doctor.}//
\endgl
\xe

\noindent In both orders, \ird{staly} remains clause-final, while the
initial \synhead{DP} supplies the topic. With \ird{Marek} as topic, the
negative denies a property ascription; with \ird{doktor} as topic, it denies
Marek's identification as the doctor.

The same predicate contrast extends beyond profession terms. In the
following schemata, the direct-case slots represent complete nominal
expressions, while the verbal slot represents a complete negative finite
predicate (§\,\ref{sec:lexical-slot-schemata}):

\pex\label{ex:neg-predicate-schemata}
\a
\begingl
\gla \textnormal{[[this-building]]}\textsubscript{\Dir{}} \textnormal{[[school]]}\textsubscript{\Dir{}} staly.//
\glb this.building.\Dir{} school.\Dir{} \Neg{}//
\glft \trsl{This building is not a school.}
\lingcontext{Correction of its classification; nominal predication.}//
\endgl
\a
\begingl
\gla \textnormal{[[this-building]]}\textsubscript{\Dir{}} \textnormal{[[be-empty]]}\textsubscript{\textsc{neg.cont}}.//
\glb this.building.\Dir{} be.empty.\Neg{}.\Cont{}//
\glft \trsl{This building is not empty.}
\lingcontext{A finite state predicate, with the regular indicative negative.}//
\endgl
\a
\begingl
\gla \textnormal{[[medicine]]}\textsubscript{\Dir{}} zada.//
\glb medicine.\Dir{} \NegExst{}-\Cont{}//
\glft \trsl{There is no medicine.}
\lingcontext{Denial of availability in the relevant situation.}//
\endgl
\xe

The verbal placeholder in (\ref{ex:neg-predicate-schemata}b) includes
\ird{za-}/\ird{z-}; its onset-dependent allomorph cannot be selected until
the lexical form is supplied. The distinction between nominal classification,
finite state predication and existential absence determines the negative
construction, whereas an English translation with \trsl{is not} does not.

Invariant \ird{staly} realises Neg\textsuperscript{0}$[u\textsc{neg}]$, which selects the
equative predication directly, without an intervening aspectual layer. The
silent operator $\mathrm{Op}_{\neg}$$[i\textsc{neg}]$ in Spec,\synhead{NegP} contributes semantic
negation, while the topic occupies the higher Spec,\synhead{TopP}:

\begin{center}
\begin{forest}
  syntax tree
  [\synhead{TopP}
    [{\synhead{DP}$_i$\\\textit{Marek}\\{\footnotesize\textsc{[dir]}}\\{\scriptsize topic}}]
    [{Top$'$}
      [\synhead{NegP}
        [{$\mathrm{Op}_{\neg}$\\{\footnotesize$[i\textsc{neg}]$}}]
        [{Neg$'$}
          [\synhead{EquP}
            [{\synhead{DP}\\$t_i$}]
            [{Eq$'$}
              [{\synhead{DP}\\\textit{doktor}\\{\footnotesize\textsc{[dir]}}\\{\scriptsize pred.}}]
              [{Eq$^{0}$\\{$\varnothing$}}]
            ]
          ]
          [{Neg$^{0}$\\\ird{staly}\\{\footnotesize$[u\textsc{neg}]$}}]
        ]
      ]
      [{Top$^{0}$\\$\varnothing$}]
    ]
  ]
\end{forest}
\end{center}

\noindent The tree omits the intervening phonologically null inflectional and
left-peripheral projections. The topic's lower copy occupies the equative
specifier, while the predicative \synhead{DP} remains the complement of
Eq\textsuperscript{0}; the inner Spec,\synhead{NegP} hosts the operator.
The participant's intermediate outer negative edge is omitted.

The distribution of \ird{staly} and \ird{za-} is conditioned by predicate
type (§\,\ref{sec:vm-neg-exponents}). Zero-copula
clauses take \ird{staly}, while the overt transitional copular verb
\ird{neh-}, being a finite indicative verb, takes \ird{za-}, giving the
negative stem \ird{za-neh-} (§\,\ref{sec:clausal-negation}).

Affirmative \ird{Marek vy doktor} retains \ird{doktor} as the nominal
 predicate within EquP: only its referential participant undergoes ordinary
 placement. Negative \ird{Marek doktor vy staly} has a different pronounced
 predicate position. The negative equative analysis assigns a participant
 escape requirement to Neg\textsuperscript{0}, realised as \ird{staly},
 and, when Phas is present, an EquP-remnant requirement to the outer edge of
 PhasP. First the equative participant moves to an outer specifier of NegP,
 above the silent negative operator; the EquP remnant then raises to the outer
 Spec,\synhead{PhasP}, preceding the phasal particle in its inner specifier.
 Finally Ord targets the participant's accessible occurrence at the negative
 edge, from which it continues to OrdP and, if independently licensed, TopP.
 The participant-escape, remnant and ordinary-placement requirements are
 distinct: the remnant occupies an outer phasal specifier, whereas tucking-in
 governs repeated satisfaction of the ordinary-placement requirement.
 Ord does not extract the participant's lower copy from the displaced EquP.
 This sequence evacuates the participant before remnant movement and
 completes each lower step before building the higher placement edge:
\[
 \begin{aligned}
 [\mathrm{EquP}\ D\ N_{\mathrm{pred}}\ \mathrm{Eq}^{0}]
 &\ \longrightarrow\ [D_i\ [\mathrm{Op}_{\neg}\ [\mathrm{EquP}\ t_i\ N_{\mathrm{pred}}\ \mathrm{Eq}^{0}]\ \mathrm{staly}]]\\
 &\ \longrightarrow\ [\mathrm{EquP}]_j\ [\mathrm{vy}\ [D_i\ [\mathrm{Op}_{\neg}\ t_j\ \mathrm{staly}]]]\\
 &\ \longrightarrow\ D_i\ [\mathrm{EquP}]_j\ [\mathrm{vy}\ [t_i\ [\mathrm{Op}_{\neg}\ t_j\ \mathrm{staly}]]].
 \end{aligned}
\]
The outer phasal edge requirement is restricted to this negative equative
 construction; it does not turn a nominal predicate into an ordinary argument
 goal of Ord. Interpretation uses the lower EquP occurrence under negation
 and Phas, retaining $\mathrm{PHAS}(\neg\mathrm{Equ})$, while the participant
 is pronounced only at its highest licensed occurrence. Without Phas,
 participant escape and ordinary placement suffice for \ird{Marek doktor staly};
 no remnant displacement above Phas is required. The remnant step temporarily
 places a participant copy above its antecedent; the later Ord step restores
 the participant's higher pronounced position. Locality is checked at each
 movement step, without imposing proper binding on every intermediate
 representation. The distinction between derivational locality and such a
 representational requirement is discussed by Müller \cite[§1]{mueller2000};
 it supplies no general permission to extract from a moved phrase.
 Thus nominal predicates
 neither move indiscriminately with arguments nor invariably follow every
 particle. Existential \ird{Pivo vy zada}, by contrast, places the nominal
 pivot before Phas while retaining the existential predicate and its negative
 suppletion in the lower predication (§\,\ref{sec:existential-constructions}).

\section{Identificational negation}\label{sec:negative-identificational-constructions}\label{sec:vm-ident-neg}
A participant nominal predicate identifies the topic with an individual
satisfying its description (§\,\ref{sec:nmlz-predicative}); clause-final
\ird{staly} denies that identification. This remains propositional negation,
although its predicate differs from a finite event or state predicate.
When the discourse treats the event and an identifiable participant as given,
the negative retains that background and commonly suggests another participant:

\pex
\a
\begingl
\gla Lobek Janku zapjaštník.//
\glb apple.\Dir{} Janek.\Trs{} \Neg{}-eat-\textsc{trs}-\Pf{}//
\glft \trsl{Janek did not eat the apple.} \lingcontext{propositional negation: denies that the event occurred}//
\endgl
\a
\begingl
\gla Lobek Jancí pjaštníkou staly.//
\glb apple.\Dir{} Janek.\Ind{} eat-\textsc{trs}-\Pf{}-\Nmlz{} \Neg{}//
\glft \trsl{The apple is not what Janek ate.} \lingcontext{identificational negation: event presupposed; identification denied}//
\endgl
\xe

\pex
\a
\begingl
\gla Janek lúbcí pjáščikou staly.//
\glb Janek.\Dir{} apple.\Ind{} eat-\Antip{}-\Pf{}-\Nmlz{} \Neg{}//
\glft \trsl{Janek is not the one who ate the apple.} \lingcontext{agent-identification denied: it was someone else}//
\endgl
\xe

\noindent The finite negative \ird{Lobek Janku zapjaštník} leaves open
whether Janek ate anything, whereas the identificational negatives reject a
participant identification against an established event. In the patient-nominal
construction, the apple is rejected as what Janek ate; in the agent-nominal
construction, Janek is rejected as the eater. In both, the suggestion of an
alternative participant depends on retaining the existential background of the
nominal description.

Existence and uniqueness requirements in contexts with no established event
or with several possible participants are undocumented. The contrast with
finite verbal negation is summarised below.

\medskip

\noindent\begin{tabularx}{\linewidth}{l L L}
\toprule
& \textsc{finite verbal negative} & \textsc{identificational negative} \\
\midrule
Predicate & finite event or state predicate & participant nominal \\
What is denied & verbal predication & identification of the topic \\
Event status & need not be given & commonly given in corrective uses \\
Negative form & \ird{za-} on the verbal complex & clause-final \ird{staly} \\
\bottomrule
\end{tabularx}

\medskip


\section{Suppletive exclusivity: \ird{jec-}, \ird{zad-}, and the \synhead{NegP} alternation}\label{sec:neg-exst-alternation}%
\label{sec:vm-jec-zad}%
\label{sec:vm-jec}

The alternation \ird{jec-}/\ird{zad-} is analysed as polarity-conditioned
suppletion at Exist\textsuperscript{0}, with a higher Neg head conditioning
\ird{zad-} and receiving a zero exponent
(§\,\ref{sec:vm-neg-exponents}). The negative
contribution is supplied by $\mathrm{Op}_{\neg}$$[i\textsc{neg}]$ in Spec,\synhead{NegP}, while
Neg\textsuperscript{0}$[u\textsc{neg}]$ conditions its overt morphological expression at
the existential root. Semantically, an affirmative existential asserts that an individual satisfies
the pivot description, while the negative existential denies that any such
individual exists. In the syntactic analysis, the existential contribution belongs
to Exist and the negative contribution to the operator in Spec,\synhead{NegP},
while Neg\textsuperscript{0}$[u\textsc{neg}]$ conditions the overt root form \ird{zad-}.
The same opposition is preserved across all
construction types --- bare existential, locative, possessive, and
presentational --- and there is no grammatical asymmetry between the two roots
in terms of which construction types each permits.

In presentational constructions, both \ird{jeca} and \ird{zada} have fixed
continuous forms. The affirmative introduces a referent satisfying an
event-anchored description, whereas the negative denies that any such referent
exists, with the description itself supplied by the discourse.

\noindent\textit{The exclusion of \ird{*zajec-}.}\enspace The existential root \ird{jec-}
does not take the indicative negative marker \ird{za-}, so that forms such as
\ird{*zajeca} and \ird{*zajecime} are excluded across the constructions
described here, which instead use \ird{zad-}:

\pex
\a
\begingl
\gla Bylam na tóm zada.//
\glb child.\Ind{} in/on book.\Dir{} \NegExst{}-\Cont{}//
\glft \trsl{The child does not have a book.} \lingcontext{\ird{zad-} obligatory in the negative possessive}//
\endgl
\a
\begingl
\gla {*}Bylam na tóm zajeca.//
\glb child.\Ind{} in/on book.\Dir{} \Neg{}-\Exst{}-\Cont{}//
\glft \trsl{(intended: the child does not have a book)} \lingcontext{\ird{za-} cannot prefix \ird{jec-}}//
\endgl
\xe

\pex
\a
\begingl
\gla Kářnikový senica zada.//
\glb read-\textsc{trs}-\Pf{}-\textsc{rel} news.\Dir{} \NegExst{}-\Cont{}//
\glft \trsl{There is no news that (I) read.} \lingcontext{\ird{zad-} obligatory in the negative presentational}//
\endgl
\a
\begingl
\gla {*}Kářnikový senica zajeca.//
\glb read-\textsc{trs}-\Pf{}-\textsc{rel} news.\Dir{} \Neg{}-\Exst{}-\Cont{}//
\glft \trsl{(intended: there is no news that (I) read)} \lingcontext{\ird{za-} cannot prefix \ird{jec-}}//
\endgl
\xe

Negative existential clauses combine \ird{zad-} at Exist\textsuperscript{0}
with zero at Neg\textsuperscript{0}, excluding the affirmative root in
\ird{*zajeca} and the additional prefix in \ird{*zazada}. These restrictions
form part of the contextual spellout of existential polarity, under which the
negative head remains syntactically present but has no overt exponent of its
own.


\section{Negative indefinites and strict concord}
\label{sec:nm-negative-concord}\label{sec:vm-negative-concord}
\ird{Deno} \trsl{nothing}, \ird{zide} \trsl{no one}, \ird{žena}
\trsl{nowhere}, \ird{žemě} \trsl{never}, and adnominal \ird{žé}
\trsl{not one} obligatorily co-occur with the appropriate realisation of
clausal negation in an ordinary finite clause, irrespective of their surface
position. Their co-occurrence expresses a single semantic negation, giving
Iridian \term{strict negative concord} in the distributional sense described
by Giannakidou \cite{giannakidou2006}. Indicative verbal predicates require
\ird{za-}, equatives require \ird{staly}, and existential predicates require
\ird{zad-}; a negative non-indicative clause instead uses its dedicated Mood
exponent (§\,\ref{sec:vm-neg-exponents}; §\,\ref{sec:vm-moodp}). A
negative indefinite cannot itself replace any of these markers.

\pex\label{ex:nm-negative-concord}
\a
\begingl
\gla Deno že zastráda.//
\glb nothing.\Dir{} 1\Sg{}.\Ind{} \Neg{}-lack-\Cont{}//
\glft \trsl{Nothing is lacking to me; I lack nothing.}//
\endgl
\a \ljudge{*}\ird{Deno že stráda.}
\lingcontext{intended: nothing is lacking to me; unlicensed negative indefinite}
\xe

\noindent In this construction, \ird{strád-} takes a direct-case theme and
an indirect-case experiencer, so that \ird{deno} supplies the theme and
\ird{že} the experiencer. The continuous \ird{-a} supplies no future tense;
the rendering \trsl{I shall not want} is available in a context that supports
a prospective interpretation. Abstracting away from aspect, the clause has
the interpretation $\neg\exists x\,[\operatorname{thing}(x)\land
\operatorname{lack}(x,\mathrm{me})]$. Initial \ird{deno} remains within the
scope of negation at interpretation in the lower argument domain
(§\,\ref{sec:tc-quantificational-pivots}).
Ordinary placement pronounces the theme above the particle field while
retaining its lower occurrence within the negative operator's scope. The
negative indefinite and the operator thus retain their clause-local concord
dependency; high pronunciation neither supplies a second negative operator
nor gives the indefinite unrestricted wide scope.

The same licensing requirement holds with the other predicate types and
with negative Mood:

\pex\label{ex:nm-concord-predication}
\a
\begingl
\gla Zide doktor staly.//
\glb no.one.\Dir{} doctor.\Dir{} \NegCop{}//
\glft \trsl{No one is a doctor.} \lingcontext{equative concord}//
\endgl
\a
\begingl
\gla Deno zada.//
\glb nothing.\Dir{} \NegExst{}-\Cont{}//
\glft \trsl{Nothing exists.} \lingcontext{existential concord}//
\endgl
\a
\begingl
\gla Deno pjašté.//
\glb nothing.\Dir{} eat-\Proh{}//
\glft \trsl{Do not eat anything.} \lingcontext{prohibitive concord; null addressee}//
\endgl
\xe

The grammar analyses these expressions as existential indefinites bearing
an uninterpretable feature $[u\textsc{neg}]$, licensed by the $[i\textsc{neg}]$ feature of a
clause-local covert negative operator. This adopts the syntactic agreement
approach developed by Penka and Zeijlstra \cite{penka2011,zeijlstra2022}, while preserving the
distinction between licensing and Vocabulary Insertion: the indefinite and
the negative predicate morphology participate in the same dependency, whose
semantic negation is supplied once by the operator. Multiple Agree permits
several $[u\textsc{neg}]$ expressions to share that licenser \cite{penka2012}. Thus the schematic sequence
\ird{zide \ldots{} deno \ldots{} za-V}, with each argument assigned its
appropriate case, expresses one negation over two existential variables;
additional concord items do not introduce further negative operators. The
existential treatment is an analysis of Iridian, rather than a definition of
all cross-linguistic n-words: Giannakidou contrasts
existential and universal treatments \cite{giannakidou2006}.

Japanese and Korean provide a useful contrast at the level of analysis.
Sells discusses Japanese \textit{dare-mo} and Korean \textit{amwu-to} with
overt predicate negation, but argues from intervention and focus that these
items can take negation in their immediate scope, rather than being
existentials interpreted below it \cite[§§\,1.2--2]{sells2022}. Iridian's
co-occurrence requirement therefore does not by itself establish its
$\neg\exists$ analysis: that interpretation is specified by the local concord
system described here. The Japanese--Korean comparison separates the
observable requirement for negative morphology from the quantificational
analysis of the dependent expression.

The licensing domain is the local finite clause. A negative operator in a
matrix clause does not license a negative indefinite inside a finite
complement (§\,\ref{sec:clausal-negation}; §\,\ref{sec:vm-po-comp}). Interrogative \ird{by}
and non-assertive complement marking do not supply $[i\textsc{neg}]$, so that questionhood
or non-assertion alone cannot license a member of the negative series.
Conditionality likewise supplies no negative licenser for an antecedent
(§\,\ref{sec:cond-negation}).
Ordinary indefinites and the universal series lack $[u\textsc{neg}]$ and do not enter
this concord dependency; the wider distribution of polarity-sensitive
indefinites is a separate descriptive question.

Short answers such as \ird{deno} or \ird{zide} are analysed as remnants of an
elliptical negative clause whose operator and predicate morphology are
recoverable. Their negative interpretation therefore does not license an
independent negative-quantifier use in a fully pronounced affirmative clause.
Genuine logical double negation requires two semantic negative operators,
ordinarily in distinct clauses or propositions; the combination of a negative
indefinite with its local predicate negative cannot produce it.

\FloatBarrier
